% ECONOMICS_PROBLEM: {"id": "D83-459", "title": "Personalization and competition", "jel_code": "D83", "source_id": "OP-0400"}
% !TeX program = xelatex
\documentclass[11pt,a4paper]{article}
\usepackage[margin=23mm]{geometry}
\usepackage{fontspec}
\setmainfont{Latin Modern Roman}
\usepackage{amsmath,amssymb,mathtools}
\usepackage{xcolor,enumitem,booktabs,longtable,array,xurl,hyperref}
\definecolor{muted}{HTML}{53636C}
\hypersetup{colorlinks=true,urlcolor=blue,linkcolor=blue}
\setlength{\parindent}{0pt}
\setlength{\parskip}{5pt}
\newcommand{\R}{\mathbb R}
\newcommand{\E}{\mathbb E}
\newcommand{\N}{\mathbb N}
\newcommand{\ind}{\mathbf 1}
\newcommand{\Var}{\operatorname{Var}}
\newcommand{\Cov}{\operatorname{Cov}}
\newcommand{\supp}{\operatorname{supp}}
\DeclareMathOperator*{\argmin}{arg\,min}
\DeclareMathOperator*{\argmax}{arg\,max}
\newcommand{\entrymeta}[3]{{\sffamily\small\color{muted}\textbf{\texttt{#1}}\quad|\quad #2\par #3\par}}
\newcommand{\Problemref}[1]{\texttt{#1}}
\newcommand{\JELref}[2]{\texttt{#2} (JEL \texttt{#1})}
\begin{document}
\section*{Personalization and competition}
\textbf{Problem ID:} \texttt{D83-459}\quad \textbf{JEL:} \texttt{D83}\par
% BEGIN ECONOMICS STATEMENT
\label{problem:OP-0400}
{\sffamily\small\color{muted}Original subject group: Digital, platform and data economics\par}
\label{definitions:problem:30007682}
\textbf{Audit classification:} Published research agenda. The source supports platform-design and equilibrium-response research. Maximum attainable utility must respect users information and search constraints rather than grant full-information choice for free.
\subsection*{Definitions and precise research target}
Consider a travel-booking platform with a fixed population of prospective users and participating accommodation sellers during a twelve-month study. A ranking policy \(r\in\{0,1\}\) is either a common ranking, \(r=0\), or personalized to observed user characteristics, \(r=1\). Let \(p_j(r)\) and \(q_j(r)\) be seller \(j\)'s equilibrium price and service quality after sellers can adjust to the announced policy. Each user \(i\) has money-valued utility \(u_i(j,q_j)-p_j-c_i\) from booking seller \(j\), where \(c_i\) is the search cost generated by the displayed ranking; the outside option is normalized to zero. Let \(CS(r)\) be expected maximum attainable utility after the user's optimal search, averaged over this user population. The target is \(\Delta CS=CS(1)-CS(0)\) after seller adjustment, together with effects for prespecified user groups. Estimate demand and search behavior using randomized rankings while accounting for competition across sellers and treatment spillovers. Combine experimental variation with a clearly stated equilibrium pricing model or market-level policy randomization to recover long-run responses. Compare this target with a counterfactual holding prices and quality fixed. The difference measures how supplier responses alter apparent personalization benefits. The specification addresses a defined marketplace and horizon, rather than treating short-run click gains as a general welfare result.

\textbf{Definitions, assumptions and mathematical qualifications.} Users' values are quasi-linear in money. A search policy \(\pi_i\) is a stopping-and-inspection rule measurable in displayed information, with total cost \(c_i(\pi_i,r)\); it cannot use unobserved seller qualities. Define \(CS(r)=E[\sup_{\pi_i}E[v_i(j(\pi_i),q(r))-p_{j(\pi_i)}(r)-c_i(\pi_i,r)\mid\text{initial information}]]\), where the no-booking action is feasible. Seller prices, quality and entry are generated by a fully specified game and equilibrium selection after policy announcement. A twelve-month endpoint is a fixed-horizon response, not an infinite-horizon stationary equilibrium unless convergence is established.
\subsection*{Variants and assumption changes}
\begin{enumerate}
\item \textbf{Editorial, immediate response.} Hold seller prices, quality and participation at baseline while randomizing rankings, and estimate the direct search-and-match benefit.
\item \textbf{Editorial, market rollout.} Randomize whole geographically isolated markets, allow suppliers to adjust for twelve months, and compare that contrast with the fixed-supplier effect; include group-specific surplus and outside-option use.
\end{enumerate}
\subsection*{References and source audit}
\textbf{Checked source.} 2025 conference draft, Sections 1 and 1.1, PDF pp. 2--5. The discussion motivates interaction, disclosure, summaries, trust, monitoring and equilibrium research. Full primary text retrieved. \url{https://conference.nber.org/confer/2025/DTs25/farronato.pdf}.
\begin{enumerate}\raggedright
\item\hypertarget{definitions:source:30007682:1}{} Chiara Farronato. 2025. \emph{Digital Platforms as Information Aggregators: Implications for their Design and Regulation}. 2025 conference draft, Sections 1 and 1.1, PDF pp. 2--5.
\url{https://conference.nber.org/confer/2025/DTs25/farronato.pdf}.
\end{enumerate}

\subsection*{Common conventions: Source mathematical conventions}

These conventions apply to each entry unless explicitly replaced.

\textbf{Models and identification.} A primitive model \(m\) specifies all probability laws, feasible choices, information, equilibrium and measurement rules used by its target. The admissible class \(\mathcal M\) is fixed independently of outcomes. If \(P_O(m)\) is its observable probability law and \(\tau(m)\) the target, its identified set at observable law \(P\) is
\[
 \mathcal I_\tau(P)=\{\tau(m):m\in\mathcal M,\ P_O(m)=P\}.
\]
Point identification means that this set is a singleton. An empty set signals incompatibility of data and assumptions. Coordinate bounds need not describe the joint set. Bounds are sharp only if every retained target is attainable or arbitrarily approachable by compatible models. Finite bounds require explicit restrictions on unobserved quantities; unrestricted confounding, hidden wealth, extrapolation or arbitrary equilibrium selection can yield uninformative sets.

\textbf{Causal interventions.} A potential outcome \(Y_i(p)\) is the outcome under a complete policy regime \(p\), including timing, financing, information and market spillovers. The target population and units are fixed before treatment. With interference, use the complete assignment vector or a justified exposure mapping; individual no-interference notation alone cannot identify an economy-wide intervention. A difference of mean potential outcomes does not require the cross-world joint distribution. Conditioning on post-treatment migration, survival, enrollment or employment changes the estimand and needs separate justification.

\textbf{Statistical statements.} An estimator, confidence set, forecast or test specifies a sampling law, model class and loss/coverage criterion. Uniform \((1-\alpha)\)-coverage means \(\inf_{m\in\mathcal M}\Pr_m\{\tau(m)\in C_n\}\geq1-\alpha\), or an explicitly asymptotic version. The whole parameter space trivially has coverage, so informativeness requires a stated width, risk or power objective. A forecasting improvement is relative to a named benchmark, information set and evaluation population.

\textbf{Equilibrium and optimization.} An equilibrium comprises feasible plans, optimality, beliefs where needed, budget constraints and all market/resource clearing. The primitives or a stated benchmark define each correspondence \(\mathcal E(p;m)\). Nonemptiness is assumed or proved, never inferred just from compact policy sets. A selected-equilibrium target uses a declared selection rule; a robust target quantifies over all permitted equilibria. Use suprema or infima unless attainment is established. An \(\varepsilon\)-optimal policy is feasible and within \(\varepsilon\) of the finite optimum value. Report an empty feasible set separately.

\textbf{Welfare and accounting.} Normative utility, interpersonal normalization, social weights, numeraire, discounting and the use of public receipts are primitives. Behavioral utility need not coincide with the welfare criterion. Household welfare includes ownership income and rents once; adding profits again to compensating variation generally double-counts them. A monetary equivalent is defined by an explicit transfer and price convention, with monotonicity and range conditions. Average effects, mechanism contributions, transfers and welfare losses are different targets. Infinite sums require convergence and dynamic feasible plans require terminal or no-Ponzi conditions.

\textbf{Source versions.} A working-paper date, revision date and journal publication year are distinguished. A DOI for a journal version is not automatically the DOI of an earlier manuscript. Locators refer to the stated version and printed pages where specified. A metadata-only check does not verify a claim about a full-text conclusion. Every entry's source subsection narrows the provenance claim accordingly.
% END ECONOMICS STATEMENT
\end{document}
